\section{はじめに}
\label{sec:intro}

野球における変化球の軌道は，リリース時の球速・回転数・回転軸の向き，および球体表面の縫い目の位置によって決まる空気力学的な力（抗力・揚力・非マグヌス力）の積分として定まる。
なかでも縫い目まわりの流れの剥離点が回転軸方向に応じて非対称に移動することで生じる非マグヌス力は，スウィーパーに代表される近年の変化球の大きな変化量を説明する要因のひとつと考えられているが，その大きさは風洞実験や実測データからの間接的な推定に頼っており，流体の詳細な挙動そのものを直接観測することは難しい。

本ノートは，格子ボルツマン法（Lattice Boltzmann Method, LBM）による3次元非定常流体解析と，ボールの並進・回転運動を記述する剛体力学とを1ステップごとに結合する数値シミュレーションにより，縫い目付き球体まわりの乱流場とそれが生成する軌道を再現することを目的とするプロジェクトの，理論的背景と実装方法をまとめたものである。

\subsection{LBMを採用する理由}

対象とする流れは，ボール直径程度の代表長さに対してレイノルズ数が\(10^5\)のオーダーに達する乱流であり，かつ球体表面の縫い目という数ミリメートル規模の幾何学的特徴が剥離点の位置を左右する。
このような問題に対しては，
\begin{itemize}
	\item 移動・回転する複雑形状境界を陽に扱いやすいこと，
	\item 各格子点の更新が局所的な演算のみで完結し，GPUによる大規模並列化と親和性が高いこと，
	\item 圧縮性ナビエ・ストークス方程式を弱圧縮性極限で近似的に解くことで，非定常な渦構造を陽に解像できること，
\end{itemize}
の3点から，格子ボルツマン法を採用した。
特にGPUによる並列計算との親和性は，1球の投球について本塁到達までの飛行時間（約0.5秒）を3次元乱流解析として計算する上で実用上の制約となる計算コストを扱える範囲に収めるために不可欠である。

\subsection{密結合方式によるボールとの相互作用}

ボールに働く力とトルクは流体解析から瞬時に決まり，その力とトルクがボールの並進・回転運動を変化させ，運動が変化した結果として境界条件が更新され，それが次の瞬間の流体解析に影響する。
本プロジェクトでは，この双方向の依存関係を近似せずに毎ステップ解くという意味で，これを「密結合方式」と呼んでいる。
具体的には，ボールを常に計算領域の中心に固定した「ボール追従座標系」を採用し，ボールの並進運動を計算領域全体に課す一様な体積力（非慣性系の見かけの力）として実現することで，領域を移動させることなく長距離の飛行を計算できるようにしている。
密結合方式の詳細な導出は\secref{coupling}で述べる。

\subsection{本ノートの構成}

本ノートは，開発の過程で得られた検証結果や設計判断を随時記録している設計ノート（本リポジトリの\texttt{docs/design/DESIGN.md}）をもとに，理論と実装の要点を整理し直したものである。構成は次のとおりである。

\begin{itemize}
	\item \secref{aerodynamics}：野球ボールの空気力学的背景（無次元パラメータ，抗力，マグヌス力・非マグヌス力，運動方程式）
	\item \secref{lbm}：格子ボルツマン法の基礎理論（ボルツマン方程式，離散速度セット，格子単位系）
	\item \secref{cumulant}：数値的安定性のための中心モーメント衝突演算子
	\item \secref{turbulence}：乱流モデリング（LES，Smagorinskyモデル，壁面近傍の扱い）
	\item \secref{geometry-boundary}：縫い目形状のパラメトリックモデリングと移動境界条件
	\item \secref{coupling}：ボールとの密結合による軌道計算アルゴリズム
	\item \secref{gpu}：GPUによる並列実装
	\item \secref{validation}：解析解・実験値との比較による検証（Validation \& Verification）
\end{itemize}
